\documentclass[10pt,a4paper]{amsart}
\usepackage[margin=19mm]{geometry}
\usepackage{amsmath,amssymb,mathtools}
\usepackage[scr=rsfs,cal=euler]{mathalfa}
\usepackage{xfrac}
\usepackage[unicode,hidelinks,bookmarks=false]{hyperref}
\newcommand{\ie}{i.e.\ }
\newcommand{\cf}{cf.\ }
\newcommand{\resp}{respectively\ }
\newcommand{\Subsection}{\subsection}
\providecommand{\nobreakdash}{\nobreak\hbox{-}}
\setlength{\emergencystretch}{2em}
\multlinegap0pt
\usepackage{amssymb}
\usepackage{tikz}
\newtheorem{lemma}{Lemma}
\newtheorem{theorem}{Theorem}
\newcommand{\Cyc}{\mathcal{C}}
\newcommand{\TC}{\mathrm{TC}}
\title{Small graphs without power-of-two cycles:\\ a lower bound of 24, a correction to a construction of Exoo,\\ and explicit bounds for $f(k)$}
\author{Daniel Garcia}
\date{Source: arXiv:2609.04686v1 (CC BY 4.0)}
\begin{document}
\maketitle

\noindent\emph{Source and licence.} Small graphs without power-of-two cycles:\\ a lower bound of 24, a correction to a construction of Exoo,\\ and explicit bounds for $f(k)$. Version arXiv:2609.04686v1 (CC BY 4.0), distributed by arXiv under the Creative Commons Attribution 4.0 International licence, http://creativecommons.org/licenses/by/4.0/, as recorded in the arXiv licence metadata for this exact version and shown on the abstract page https://arxiv.org/abs/2609.04686v1. Full licence notice: \texttt{licenses/arxiv-2609.04686-CC-BY-4.0.txt}.\par\smallskip

\noindent\textit{Editorial scope.} Counted unit: the article's theorem thm:450 (source0.tex lines 184--186). The article's complete original proof of it is delivered with it (source0.tex lines 188--190, one \texttt{proof} environment of 3 lines). No mathematics is written, repaired or re-derived: the statement and the proof are the article's own lines, in the article's own notation, and no retained line is substituted or re-flowed.  Retained uncounted as the source-local context the proof needs: lines 121--123.  Nothing was omitted from any retained span, so the package carries no omission record.  No retained line contains a citation, so the wrapper carries no bibliography and no bibliography entry.  No retained line contains a figure environment, an image-inclusion command or drawing code, so no image is delivered and the package declares no asset dependency.  Editorial formatting only, in addition: an AMS \texttt{amsart} preamble that restates the article's own declarations and macros used by the retained lines, namely the environments \texttt{lemma}, \texttt{theorem} on separate counters, together with the standard packages those lines need.  The retained environments print in this document as Lemma 1, Theorem 1 in order of appearance, whereas the article prints the same items with the numbering of its own full document.  The complete native source of the article as distributed in the arXiv e-print for this exact version is bundled unchanged as \texttt{sources/209381/source0.tex}.
\begin{lemma}[Window lemma]\label{lem:window}
Let $B$ be cubic and let $G$ be obtained by replacing every vertex $x$ of $B$ by a gadget $H_x$. Then a cycle of $G$ either lies inside a single gadget, or projects onto a cycle $x_1\cdots x_\ell$ of $B$ and has length $\ell+\sum_{i=1}^{\ell}s_i$ with $s_i\in S_{H_{x_i}}(p_i,q_i)$, where $p_i,q_i$ are the attachments facing $x_{i-1}$ and $x_{i+1}$. Conversely, every such sum is the length of a cycle of $G$. In particular $G$ has a cycle of length $T$ if and only if some gadget has an internal cycle of length $T$ or $T-\ell$ lies in the Minkowski sum $\sum_i S_{H_{x_i}}(p_i,q_i)$ for some cycle of $B$ of length $\ell$.
\end{lemma}

\begin{theorem}\label{thm:450}
There is an orientation of the $H_{15}$-replacement of $\TC$ that yields a cubic graph $G'_{450}$ on $450$ vertices with no cycle of length $4$, $8$, $16$ or $32$. Hence $f(5)\le 450$.
\end{theorem}

\begin{proof}
Choose for each vertex of $\TC$ which incident edge faces $u$ so that no 8-cycle of $\TC$ has all eight of its vertices' $u$-edges on the cycle. This is a satisfiable SAT instance with 90 clauses (one per 8-cycle); a solution is listed in the appendix. For the resulting graph, Lemma~\ref{lem:window} gives: a base 8-cycle has at least one $v$--$w$ visit, so its expansions have length at least $8+7\cdot 3+5=34$; every other base cycle has $\ell\ge 10$ ($\TC$ is bipartite of girth 8) and expands to length at least $10\cdot(1+3)=40$; and $\Cyc(H_{15})$ contains no power of two. Thus no cycle of length $4,8,16$ or $32$ exists. This was confirmed computationally: exhaustive search finds no cycle of length 4, 8 or 16, and a SAT encoding of ``a cycle of length exactly 32'' is unsatisfiable.
\end{proof}

\end{document}
